% Incorporación ampliativa: CG-006, SR-032--SR-035 del inventario REV05.
% Residencia: inmediatamente después de la construcción de R36, antes de G9.
% Propietario computacional: 16_CIERRE_GLOBAL_HMT_MD_2026-07-22/
% 01_SELECTOR_CONSTANTES/selector_interno_constantes.py, líneas 31--63,
% 131--180, 185--240 y 267--321.
% Recuperación y cambio de representación: output/REVISION_EDITORIAL_HMT_MD_
% 20260905/03_PRUEBAS/verificar_generacion_infinita_nonadica.py,
% líneas 54--77, 651--723 y 974--1147.
% Exposición de origen: ampliaciones_sucesoras_20260824/parte_i_ii/owners/
% U008_orbitas_elevacion_r36.tex, líneas 210--358.
% Estatuto: RESULTADO_RECUPERADO; exposición de enumeración y naturalidad
% reunida sin promover la concordancia en R36 a equivalencia global.

\section[Sélection et transport de la première frontière]{Sélection combinatoire et transport de représentation de la première frontière conjointe}
\label{r36c:sec:frontera}

La première frontière conjointe reçoit les blocs déjà produits par les graines APP, la réduction tritique et le calendrier de relèvement du TPK. Sa construction conserve la distinction entre trois opérations : déterminer les distributions admissibles, orienter leurs deux canaux latéraux et transporter l’incidence vers le système de coordonnées exceptionnel. Cet exposé développe ces opérations sur le même état fini. Les noms des coordonnées réelles désignent leurs lecteurs ultérieurs ; les conditions de sélection qui suivent agissent sur des mots ternaires, des matrices de transport et des incidences entières.

\subsection{État régional antérieur à la frontière}

Soit $\mathcal W\subset\mathbb F_3^6$ le répertoire de $243$ mots visibles obtenu par le dénombrement des émissions APP--TRIT--TPK. Pour un vecteur ligne $s\in\mathcal W$, définissons
\[
 b_0(s)=s,\quad b_1(s)=sL_0,\quad b_2(s)=sL_0^2,\quad
 b_3(s)=sL_0^2L_1,\quad b_4(s)=sL_0^2L_1^2,
\]
avec des opérations dans $\mathbb F_3$. Les matrices précédemment dérivées des registres régionaux sont
\[
L_0=\begin{pmatrix}
2&2&2&1&2&1\\2&1&2&2&1&1\\1&0&0&1&0&2\\
0&0&1&1&1&0\\2&2&2&0&2&1\\2&1&2&1&0&0
\end{pmatrix},\qquad
L_1=\begin{pmatrix}
2&1&2&0&2&2\\1&2&1&1&2&0\\1&2&1&1&1&2\\
0&1&0&0&2&1\\2&0&2&1&0&1\\0&2&2&2&1&1
\end{pmatrix}.
\]
Le sélecteur régional antérieur fournit
\[
s_{\rm cl}=010211,\qquad s_{\rm pr}=201101,
\]
au moyen des conditions $b_2(s_{\rm cl})=b_0(s_{\rm cl})$, $b_4(s_{\rm pr})=b_2(s_{\rm pr})$ et $b_1(s_{\rm cl})=b_3(s_{\rm pr})$. À cette étape, les blocs courants sont
\[
b_4(s_{\rm cl})=110221,\qquad
b_4(s_{\rm pr})=102011.
\]
La concaténation conserve les cinq blocs de chaque registre ; le vecteur $b_4$ est l’entrée à six coordonnées reçue par la relation de frontière.

Pour exprimer fidèlement la relation récupérée, on utilise la matrice d’incidence dans sa représentation d’origine,
\begin{equation}
A^{\rm o}=\begin{pmatrix}
0&1&1&1&1&1\\1&0&1&1&2&2\\1&1&0&2&1&2\\
1&1&2&0&2&1\\1&2&1&2&0&1\\1&2&2&1&1&0
\end{pmatrix}.
\label{r36c:eq:ao}
\end{equation}
On vérifie $(A^{\rm o})^2=-I$ dans $M_6(\mathbb F_3)$ ; donc $(A^{\rm o})^{-1}=-A^{\rm o}$. Cette représentation et ses vecteurs seront ensuite transportés vers celle utilisée par le lecteur exceptionnel.

\subsection{Équation de compatibilité et saturation des colonnes}

L’axe de la frontière est fixé par $z=b_4(s_{\rm pr})=102011$. Pour chaque $s\in\mathcal W$, on définit
\begin{align}
m(s)&=b_4(s)A^{\rm o}L_1^3,\\
q(s)&=m(s)+z,\\
\widetilde q(s)&=s(A^{\rm o})^{-1}L_0^3(A^{\rm o})^3.
\label{r36c:eq:compat-datos}
\end{align}
La condition de compatibilité est $q(s)=\widetilde q(s)$. Les lignes latérales $x,y\in\mathbb F_3^6$ satisfont $x+y=m(s)$ ; la frontière ordonnée est $B=(x,y,z)^{\mathsf T}$.

Pour former le reste et la retenue, on utilise les représentants entiers $0,1,2$ de chaque trit. Si $S_j=x_j+y_j+z_j$, on exige
\[
q_j=S_j\bmod3,\qquad c_j=\frac{S_j-q_j}{3}=1,
\]
et l’occupation des colonnes est donnée par
\begin{equation}
\mathbf1_{\{x_j\ne0\}}+\mathbf1_{\{y_j\ne0\}}+\mathbf1_{\{z_j\ne0\}}
=2+\mathbf1_{\{z_j\ne0,\ m_j(s)\ne2\}}.
\label{r36c:eq:ocupacion}
\end{equation}
La neutralité duale complète la condition par
\begin{equation}
BA^{\rm o}\mathbf1=0\quad\hbox{dans }\mathbb F_3^3.
\label{r36c:eq:neutralidad}
\end{equation}
Le reste, la retenue et l’occupation interviennent conjointement : la réduction modulaire conserve le reste et le registre entier distingue les distributions qui le produisent.

\begin{lemma}[Résolution de la compatibilité régionale]
\label{r36c:lem:compatibilidad} L’équation $q(s)=\widetilde q(s)$ possède, dans $\mathbb F_3^6$, les trois solutions
\[
020010,\qquad121200,\qquad222120.
\]
Son intersection avec le répertoire visible est $\{121200,222120\}$.
\end{lemma}
\begin{proof}
L’équation s’écrit $sD=z$, où
\[
D=(A^{\rm o})^{-1}L_0^3(A^{\rm o})^3
   -L_0^2L_1^2A^{\rm o}L_1^3
 =\begin{pmatrix}
2&1&1&2&2&2\\0&2&0&0&0&1\\0&0&0&0&1&0\\
1&2&2&2&2&0\\1&2&2&0&1&2\\1&2&0&2&2&2
\end{pmatrix}.
\]
L’élimination sur $\mathbb F_3$ donne la paramétrisation
\[
s=(t,2,t,2t,1+2t,0),\qquad t\in\mathbb F_3.
\]
La substitution vérifie les six équations et en épuise les solutions. Dans le répertoire engendré $\mathcal W$, les deux mots correspondant à $t=1,2$ sont présents et celui correspondant à $t=0$ est absent. Cette dernière vérification est une interrogation exacte du dénombrement des mots visibles, antérieure à toute lecture archimédienne.
\end{proof}

\begin{proposition}[Détermination de l’axe d’auto-échelle]
La compatibilité, la retenue unitaire et l’occupation des colonnes sélectionnent $s_{\rm au}=121200$ parmi les deux solutions visibles du lemme~\ref{r36c:lem:compatibilidad}. On obtient
\[
b_4(s_{\rm au})=010200,\qquad
m(s_{\rm au})=210112,\qquad q=012120,
\]
avec l’occupation $223232$.
\end{proposition}
\begin{proof}
Pour $s=222120$, on calcule $b_4(s)=222000$, $m(s)=201010$ et $q(s)=000021$. Dans sa troisième colonne, $m_3=1$ et $z_3=2$. L’occupation requise est trois, donc $x_3,y_3$ doivent être tous deux non nuls. L’unique couple avec $x_3+y_3\equiv1\pmod3$ est $(2,2)$ ; sa somme entière avec $z_3=2$ vaut six et produit une retenue de deux. Cela contredit $c_3=1$.

Pour $s=121200$, les six colonnes admettent respectivement les couples
\[
\begin{array}{c|cccccc}
j&1&2&3&4&5&6\\\hline
(x_j,y_j)&(0,2),(2,0)&(2,2)&(1,2),(2,1)&(2,2)&(2,2)&(0,2),(2,0).
\end{array}
\]
Chaque couple satisfait les équations de reste, de retenue et d’occupation. Il existe donc huit dispositions avant l’application de la neutralité duale.
\end{proof}

\begin{theorem}[Énumération exacte des orientations de frontière]
\label{r36c:thm:dos} Les conditions précédentes déterminent exactement
\[
B_- =\begin{pmatrix}021222\\222220\\102011\end{pmatrix},
\qquad
B_+ =\begin{pmatrix}222220\\021222\\102011\end{pmatrix}.
\]
Toutes deux ont le reste $012120$, la retenue $111111$ et l’occupation $223232$.
\end{theorem}
\begin{proof}
La démonstration précédente réduit l’énumération à huit possibilités. Puisque $A^{\rm o}\mathbf1=(2,1,1,1,1,1)^{\mathsf T}$, la neutralité d’une ligne $v$ équivaut à $2v_1+v_2+\cdots+v_6=0$ dans $\mathbb F_3$. Son évaluation dans les huit dispositions est
\[
\begin{array}{c|c|c}
x&y&(xA^{\rm o}\mathbf1,yA^{\rm o}\mathbf1,zA^{\rm o}\mathbf1)\\\hline
021220&222222&(1,2,0)\\
021222&222220&(0,0,0)\\
022220&221222&(2,1,0)\\
022222&221220&(1,2,0)\\
221220&022222&(2,1,0)\\
221222&022220&(1,2,0)\\
222220&021222&(0,0,0)\\
222222&021220&(2,1,0).
\end{array}
\]
Les deuxième et septième lignes sont les seules solutions. Les sommes entières des colonnes de chacune sont $(3,4,5,4,5,3)$ et restituent le reste et la retenue déclarés.
\end{proof}

\subsection{Longueur dirigée et orientation de la frontière}

Le transport régional définit un graphe orienté sur $\mathbb F_3^6$, d’arêtes $v\to vL_0$ et $v\to vL_1$. Pour des vecteurs $u,v$ mutuellement accessibles, soit
\[
d(u,v)=\min\{n\ge0:\exists\epsilon_1,\ldots,\epsilon_n\in\{0,1\},
\ v=uL_{\epsilon_1}\cdots L_{\epsilon_n}\}.
\]
Cette longueur est une donnée combinatoire du transport. Sa somme sur les trois lignes conserve les étiquettes régionales :
\[
\mathcal L(B)=d(110221,B_1)+d(102011,B_2)+d(010200,B_3).
\]

\begin{proposition}[Sélection par longueur dirigée]
\label{r36c:prop:minimo} On a $\mathcal L(B_-)=23$ et $\mathcal L(B_+)=20$. Dans le domaine des deux frontières admissibles, le minimum est unique et sélectionne $B_+$.
\end{proposition}
\begin{proof}
Les chemins minimaux et leurs longueurs sont
\[
\begin{array}{c|c|c|c|c}
 &\text{origine}&\text{destination}&\text{mot de transport}&d\\\hline
B_-&110221&021222&01100011&8\\
 &102011&222220&10111111&8\\
 &010200&102011&1010101&7\\\hline
B_+&110221&222220&1111&4\\
 &102011&021222&001100101&9\\
 &010200&102011&1010101&7.
\end{array}
\]
Chaque mot se vérifie par multiplication successive de ses matrices. Pour vérifier aussi la minimalité, on forme les couches disjointes
\[
F_0(u)=\{u\},\qquad
F_{n+1}(u)=\bigl(F_n(u)L_0\cup F_n(u)L_1\bigr)
\setminus\bigcup_{j=0}^{n}F_j(u).
\]
Une récurrence sur $n$ démontre que $F_n(u)$ est exactement l’ensemble des sommets à distance $n$ : tout mot minimal de longueur $n+1$ provient d’un mot minimal de longueur $n$, et le quotient par les couches précédentes élimine les parcours plus courts. L’énumération obtenue avec les deux matrices présentées ci-dessus donne
\[
\begin{array}{c|rrrrrrrrrr}
u\backslash n&0&1&2&3&4&5&6&7&8&9\\\hline
110221&1&2&4&8&16&32&60&105&156&177\\
102011&1&2&3&5&10&20&37&71&119&171\\
010200&1&2&4&8&15&30&56&100&148&172.
\end{array}
\]
Les destinations du premier tableau apparaissent pour la première fois dans les couches qui y sont indiquées. Cette énumération est finie, utilise uniquement les produits modulo trois et coïncide avec le parcours en largeur du certificat exécutable. La somme des trois longueurs donne $8+8+7=23$ et $4+9+7=20$.
\end{proof}

L’orientation de la représentation APP au moyen de cylindres semi-ouverts sélectionne préalablement la même matrice $B_+$ à cette frontière. Le théorème~\ref{r36c:thm:dos} et la proposition~\ref{r36c:prop:minimo} établissent ainsi la concordance des deux sélections sur l’ensemble concret $\{B_-,B_+\}$. Le domaine de cette concordance est la première frontière conjointe. Le prolongement coinductif conserve ses propres relations d’admissibilité, de mémoire, de survie et d’orientation à chaque profondeur.

\subsection{Naturalité du changement de représentation exceptionnel}

La matrice utilisée par le lecteur exceptionnel dans la représentation active est
\[
A^{\rm a}=\begin{pmatrix}
0&1&1&1&1&1\\1&0&1&2&2&1\\1&1&0&1&2&2\\
1&2&1&0&1&2\\1&2&2&1&0&1\\1&1&2&2&1&0
\end{pmatrix}.
\]
La permutation des coordonnées
\[
R_p(v_1,v_2,v_3,v_4,v_5,v_6)
=(v_1,v_2,v_3,v_5,v_6,v_4)
\]
transporte simultanément le mot et la matrice d’incidence.

% REMATE_LOCAL_R36
\Needspace{36\baselineskip}
\begin{theorem}[Entrelacement de l’incidence de frontière]
\label{r36c:thm:cambio}
Pour tout $v\in\mathbb F_3^6$,
\[
R_p(vA^{\rm o})=R_p(v)A^{\rm a}.
\]
La transformation est inversible et conserve la neutralité duale, les restes, les retenues et les occupations, exprimés dans les coordonnées transportées. En particulier, le triplet sélectionné possède les deux représentations
\[
\begin{array}{c|c|c}
&\text{représentation d'origine}&\text{représentation exceptionnelle}\\\hline
\text{clôture}&222220&222202\\
\text{propagation}&021222&021222\\
\text{auto-échelle}&102011&102110.
\end{array}
\]
\end{theorem}
\begin{proof}
Si $p=(1,2,3,5,6,4)$, la comparaison des entrées des deux matrices donne $A^{\rm a}_{ij}=A^{\rm o}_{p(i),p(j)}$. Pour la coordonnée $j$,
\[
\bigl(R_p(v)A^{\rm a}\bigr)_j
=\sum_i v_{p(i)}A^{\rm o}_{p(i),p(j)}
=\sum_kv_kA^{\rm o}_{k,p(j)}
=\bigl(R_p(vA^{\rm o})\bigr)_j.
\]
L’inverse est $R_p^{-1}(w)=(w_1,w_2,w_3,w_6,w_4,w_5)$. Comme $R_p$ permute les coordonnées, il conserve leur somme et transporte coordonnée par coordonnée les sommes entières, leurs restes modulo trois, leurs quotients et les nombres d’entrées non nulles. Il en résulte la conservation de tous les registres indiqués. Le tableau s’obtient en appliquant $R_p$ aux trois mots de $B_+$, et l’inverse récupère exactement leurs mots originaux.
\end{proof}

Le lecteur positionnel et le lecteur d’incidence demeurent associés au même état enrichi par ce changement explicite de représentation. Si $\ell_{\rm pos}$ est le lecteur dans la représentation d’origine, son expression dans la représentation transportée est $\ell_{\rm pos}\circ R_p^{-1}$. Cette composition conserve l’ordre positionnel : changer la représentation de l’incidence et changer le lecteur sont des opérations coordonnées. La première frontière entre ainsi dans le prolongement avec ses trois mots, son orientation régionale, sa retenue entière et son repère exceptionnel récupérables.
